\documentclass[10pt, titlepage]{article}

\usepackage{amsmath, mathtools}

\usepackage{amsfonts}
\usepackage{amssymb}
\usepackage{graphicx}
\usepackage{colortbl}
\usepackage{xr}
\usepackage{hyperref}
\usepackage{longtable}
\usepackage{xfrac}
\usepackage{tabularx}
\usepackage{float}
\usepackage{siunitx}
\usepackage{booktabs}
\usepackage{multirow}
\usepackage[section]{placeins}
\usepackage{caption}
\usepackage{fullpage}
\usepackage{parskip}

\setlength{\parindent}{0pt}

\hypersetup{
bookmarks=true,     % show bookmarks bar?
colorlinks=true,       % false: boxed links; true: colored links
linkcolor=red,          % color of internal links (change box color with linkbordercolor)
citecolor=blue,      % color of links to bibliography
filecolor=magenta,  % color of file links
urlcolor=cyan          % color of external links
}

\usepackage{array}

\externaldocument{../../SRS/SRS}

\input{../../Comments.text}
\input{../../Common.text}

\begin{document}

\title{Module Interface Specification for \progname{}}

\author{\authname}

\date{\today}

\maketitle

\pagenumbering{roman}

\section{Revision History}

\begin{table}[!htbp]
\caption{Revision History} \label{TblRevisionHistory}
\begin{tabularx}{\textwidth}{llX}
\toprule
\textbf{Date} & \textbf{Developer(s)} & \textbf{Change}\\
\midrule
2026/11/14 & Mohammad Bilal & Added table in section 5 and completed individual reflection \\
\bottomrule
\end{tabularx}
\end{table}

\newpage

\tableofcontents

\pagenumbering{arabic}

\newpage

\section{Symbols, Abbreviations and Acronyms}

See SRS Documentation at \wss{give url}

\wss{Also add any additional symbols, abbreviations or acronyms}

\newpage

\section{Introduction}

The following document details the Module Interface Specifications for
\wss{Fill in your project name and description}

Complementary documents include the System Requirement Specifications
and Module Guide.  The full documentation and implementation can be
found at \url{...}.  \wss{provide the url for your repo}

\newpage
\section{Notation}

\wss{You should describe your notation.  You can use what is below as
  a starting point.}

The structure of the MIS for modules comes from \cite{HoffmanAndStrooper1995},
with the addition that template modules have been adapted from
\cite{GhezziEtAl2003}.  The mathematical notation comes from Chapter 3 of
\cite{HoffmanAndStrooper1995}.  For instance, the symbol := is used for a
multiple assignment statement and conditional rules follow the form $(c_1
\Rightarrow r_1 | c_2 \Rightarrow r_2 | ... | c_n \Rightarrow r_n )$.

The following table summarizes the primitive data types used by \progname. 

\begin{center}
\renewcommand{\arraystretch}{1.2}
\noindent 
\begin{tabular}{l l p{7.5cm}} 
\toprule 
\textbf{Data Type} & \textbf{Notation} & \textbf{Description}\\ 
\midrule
character & char & a single symbol or digit\\
integer & $\mathbb{Z}$ & a number without a fractional component in (-$\infty$, $\infty$) \\
natural number & $\mathbb{N}$ & a number without a fractional component in [1, $\infty$) \\
real & $\mathbb{R}$ & any number in (-$\infty$, $\infty$)\\
\bottomrule
\end{tabular} 
\end{center}

\noindent
The specification of \progname \ uses some derived data types: sequences, strings, and
tuples. Sequences are lists filled with elements of the same data type. Strings
are sequences of characters. Tuples contain a list of values, potentially of
different types. In addition, \progname \ uses functions, which
are defined by the data types of their inputs and outputs. Local functions are
described by giving their type signature followed by their specification.

\newpage
\section{Module Decomposition}

The following table is taken directly from the Module Guide document for this project \cite{MG}.

\begin{table}[H]
\caption{Module Hierarchy}
\label{TblMH}
\renewcommand{\arraystretch}{1.3}
\begin{tabularx}{\textwidth}{@{}p{0.31\textwidth}X@{}}
\toprule
\textbf{Level 1} & \textbf{Level 2} \\
\midrule
Hardware-Hiding Module & M1 Hardware-Hiding Module \\
\midrule
\multirow{18}{0.31\textwidth}{Behaviour-Hiding Module}
& M2 Application Control Module \\
& M3 Journey Progress Module \\
& M4 Wardrobe Module \\
& M5 Player Data Module \\
& M6 Preferences Module \\
& M7 Level Catalogue Module \\
& M8 Accessory Catalogue Module \\
& M9 Procedure Content Module \\
& M10 Procedure Runner Module \\
& M11 Narration Module \\
& M12 Player Select Screen Module \\
& M13 Journey Map Screen Module \\
& M14 Reward Screen Module \\
& M15 Scoreboard Screen Module \\
& M16 Dress-Up Screen Module \\
& M17 Parental Gate Module \\
& M18 Settings Screen Module \\
& M19 Procedure Guide Screen Module \\
\midrule
\multirow{7}{0.31\textwidth}{Software Decision Module}
& M20 Data File Module \\
& M21 Save Slots Module \\
& M22 Scene Staging Module \\
& M23 Character Rig Module \\
& M24 Screen Effects Module \\
& M25 Sound Effects Module \\
& M26 UI Widgets Module \\
\bottomrule
\end{tabularx}
\end{table}

\newpage
\section{MIS of \wss{Module Name}} \label{Module} \wss{Use labels for
  cross-referencing}

\wss{You can reference SRS labels, such as R\ref{R_Inputs}.}

\wss{It is also possible to use \LaTeX for hypperlinks to external documents.}

\subsection{Module}

\wss{Short name for the module}

\subsection{Uses}


\subsection{Syntax}

\subsubsection{Exported Constants}

\subsubsection{Exported Access Programs}

\begin{center}
\begin{tabular}{p{2cm} p{4cm} p{4cm} p{2cm}}
\hline
\textbf{Name} & \textbf{In} & \textbf{Out} & \textbf{Exceptions} \\
\hline
\wss{accessProg} & - & - & - \\
\hline
\end{tabular}
\end{center}

\subsection{Semantics}

\subsubsection{State Variables}

\wss{Not all modules will have state variables.  State variables give the module
  a memory.}

\subsubsection{Environment Variables}

\wss{This section is not necessary for all modules.  Its purpose is to capture
  when the module has external interaction with the environment, such as for a
  device driver, screen interface, keyboard, file, etc.}

\subsubsection{Assumptions}

\wss{Try to minimize assumptions and anticipate programmer errors via
  exceptions, but for practical purposes assumptions are sometimes appropriate.}

\subsubsection{Access Routine Semantics}

\noindent \wss{accessProg}():
\begin{itemize}
\item transition: \wss{if appropriate} 
\item output: \wss{if appropriate} 
\item exception: \wss{if appropriate} 
\end{itemize}

\wss{A module without environment variables or state variables is unlikely to
  have a state transition.  In this case a state transition can only occur if
  the module is changing the state of another module.}

\wss{Modules rarely have both a transition and an output.  In most cases you
  will have one or the other.}

\subsubsection{Local Functions}

\wss{As appropriate} \wss{These functions are for the purpose of specification.
  They are not necessarily something that is going to be implemented
  explicitly.  Even if they are implemented, they are not exported; they only
  have local scope.}

\newpage

\newpage

\newpage{}

\section{Appendix --- Generative AI Use Disclosure}

\wss{ChatGPT (version 5) was used to brainstorm ideas for the template for this disclose section.}

\subsection{AI Tools Used}

\wss{Give the name of the AI tool(s) used and the version.}

\subsection{How and Where Used}

\wss{Explain what the tool(s) were used for.  Examples include: brainstorming,
explaining a technical concept, reviewing the document, generating or modifying
text, generating or modifying diagrams, identifying errors or omissions or
inconsistencies.}

\wss{For each of the uses, identify which parts of the document were affected.}

\wss{You do not need to report every interaction, but you should disclose the 
uses that materially contributed to your work}

\subsection{Verification of AI-Generated Content}

\wss{\begin{itemize}
    \item Describe how the team checked AI-generated or AI-assisted material
    before including it in the document.
    \item In particular, verify factual claims, technical assertions,
    requirements, calculations, references, and other information for which an
    AI system may produce incorrect or fabricated results.
    \item \textbf{The team is responsible for the accuracy, completeness,
consistency, and appropriateness of the submitted document, regardless of
whether material was generated by a human or an AI system.}
\end{itemize}}

\newpage{}

\section{Appendix --- Reflection}\label{sec:Reflection}

Questions 1 and 2 are answered individually by each team member, while
Questions 3 through 7 are answered as a team. The team responses are
provided under Team Reflection, while the individual responses are provided
under each respective team member's name.

\subsubsection*{Team Reflection}

\begin{enumerate}
    \setcounter{enumi}{2}

    \item \textbf{Which of your design decisions stemmed from speaking to your
    client(s) or a proxy (e.g., your peers, stakeholders, potential users)?
    For those that were not, why, and where did they come from?}

    % Team response to be added.

    \item \textbf{While creating the design doc, what parts of your other
    documents (e.g., requirements, hazard analysis, etc.), if any, needed to
    be changed, and why?}

    % Team response to be added.

    \item \textbf{What are the limitations of your solution? Put another way,
    given unlimited resources, what could you do to make the project better?}

    % Team response to be added.

    \item \textbf{Give a brief overview of other design solutions you
    considered. What are the benefits and tradeoffs of those other designs
    compared with the chosen design? From all the potential options, why did
    you select the documented design?}

    % Team response to be added.

    \item \textbf{What did
    you learn by implementing another team's module? Were all the details you
    needed in the documentation, or did you need to make assumptions, or ask
    the other team questions? If your team also had another team implement one
    of your modules, what was this experience like? Are there things in your
    documentation you could have changed to make the process go more smoothly
    for when an ``outsider'' completes some of the implementation?}

    To be completed after the inter-team module implementation activity.

\end{enumerate}

\subsubsection*{Mohammad Bilal}

\begin{enumerate}
    \item \textbf{What went well while writing this deliverable?}

    One thing that went well was how we organized the work early in the
    process. We had to prepare two design documents while also working on
    the proof of concept (PoC) and the cognitive walkthrough. Rather than
    having everyone work on every part, we assigned one person to the
    cognitive walkthrough, two people to start the PoC, and the remaining
    two, including me, to focus on the design documents. This allowed us
    to make progress on different parts of the deliverable in parallel
    without spreading everyone too thin.

    Another thing that helped was starting the PoC early. Since we were
    all new to Unity, we did not initially have a clear understanding of
    how to organize a Unity application into the modules we needed to
    describe in the design documents. The teammates working on the PoC
    were able to explore Unity and explain what they learned to the rest
    of the team. This gave us a better understanding of which module
    boundaries made sense and how the design could actually be
    implemented, rather than designing everything without considering
    the tools we would be using.

    \item \textbf{What pain points did you experience during this
    deliverable, and how did you resolve them?}

    One challenge was the limited time available. We had to prepare two
    fairly substantial design documents, even though the design was
    still preliminary, while also completing the other parts of the
    deliverable. Trying to have everyone contribute to everything would
    have made it harder to finish on time. We addressed this by dividing
    the work into clear responsibilities, as described above, and sharing
    what we learned as we progressed. This helped keep the workload
    manageable and allowed us to work on multiple tasks at once.

    Another challenge I experienced was keeping the design documents
    consistent with the SRS and our other documentation. At this point,
    we have a large amount of material across several documents, and it
    is difficult to remember every requirement and design decision. For
    example, changes to a module's responsibilities could affect how it
    relates to requirements or other modules. I used AI to help compare
    sections and identify possible inconsistencies or missing
    connections. I then reviewed those suggestions against the SRS, the
    PoC, and our intended design before making changes. This made it
    easier to check that the modules addressed the requirements without
    introducing contradictions between documents.

\end{enumerate}

\subsubsection*{Ahsan Muzammil}

\begin{enumerate}
\item \textbf{What went well while writing this deliverable?}

One thing that went particularly well was how quickly our team
organized and divided the work. We created GitHub Issues early
on and assigned responsibilities based on each team member's
interests and strengths. Unlike previous deliverables, where
everyone worked on different sections of the same document,
we decided to divide responsibilities across the deliverables.
One member focused on the Custom Document, two worked on
the PoC, and we assigned pairs to the MIS and MG. Since we
only have five team members, those working on the PoC also
contributed to sections of the MIS and MG.

I believe this approach helped us work more efficiently
and allowed everyone to focus on areas they were most
comfortable with. Additionally, the tutorial sessions were
very helpful in clearing up confusion about the Design
Document requirements and expectations for the PoC.


\item \textbf{What pain points did you experience during this deliverable,
and how did you resolve them?}

One of the biggest challenges was the limited time available
to complete the Design Document. Given the amount of detail
required for the MIS and MG, I felt that the timeline was
quite demanding and made it difficult to spend as much time
refining the design as I would have liked.

Another major challenge for me personally was determining
how to decompose the game into appropriate software modules.
Although I had previous experience working with C\# during
my 16-month co-op, I had never worked with Unity or game
development before. As a result, understanding how to
structure a game and identify the responsibilities and
dependencies of different modules involved a significant
learning curve.

To address these challenges, we divided the workload early,
collaborated on design decisions, and used the tutorial
sessions to clarify areas of uncertainty. Reviewing the
PoC also helped me better understand how the different
parts of the game interact, which made it easier to
develop the module decomposition. While there is still
more for me to learn about Unity architecture, this
deliverable helped improve my understanding of software
design in the context of game development.

\end{enumerate}

\subsubsection*{Jackson Beach}

\begin{enumerate}
\item \textbf{What went well while writing this deliverable?}

    % Individual response to be added.

\item \textbf{What pain points did you experience during this deliverable,
    and how did you resolve them?}

    % Individual response to be added.
\end{enumerate}

\subsubsection*{Jake Finlay}

\begin{enumerate}
\item \textbf{What went well while writing this deliverable?}

    % Individual response to be added.

\item \textbf{What pain points did you experience during this deliverable,
    and how did you resolve them?}

    % Individual response to be added.
\end{enumerate}

\subsubsection*{Ethan McMehen}

\begin{enumerate}
\item \textbf{What went well while writing this deliverable?}

    % Individual response to be added.

\item \textbf{What pain points did you experience during this deliverable,
    and how did you resolve them?}

    % Individual response to be added.
\end{enumerate}

\newpage
\bibliographystyle {unsrt}
\bibliography {../../../refs/References}

\end{document}